%!TEX root = ../main.tex
% SÍMBOLOS matemáticos o de notación usados en el documento.
%
% Definir:   \newglossaryentry{sym:clave}{type=symbols, name={\ensuremath{...}},
%                description={...}, sort={clave}}
% Usar:      \gls{sym:clave}, o \glsadd{sym:clave} junto a la ecuación
%            cuando el símbolo se escribe directamente en modo matemático.
% Si no se usa ningún símbolo, la sección no se imprime.
% El campo sort fija el orden de la lista: primero la evaluación
% financiera (sección 5.x) y después el precio del resumen automático.

% --- evaluación financiera (ecuaciones del VAN, la TIR y el repago) ---
\newglossaryentry{sym:van}{
    type=symbols,
    name={\ensuremath{\mathit{VAN}}},
    description={Valor actual neto del flujo de fondos proyectado},
    sort={01}
}

\newglossaryentry{sym:tir}{
    type=symbols,
    name={\ensuremath{\mathit{TIR}}},
    description={Tasa interna de retorno: tasa que hace nulo el \ensuremath{\mathit{VAN}}},
    sort={02}
}

\newglossaryentry{sym:i0}{
    type=symbols,
    name={\ensuremath{I_0}},
    description={Inversión inicial del proyecto (período 0)},
    sort={03}
}

\newglossaryentry{sym:fc}{
    type=symbols,
    name={\ensuremath{FC_t}},
    description={Flujo de fondos neto del período \ensuremath{t}},
    sort={04}
}

\newglossaryentry{sym:r}{
    type=symbols,
    name={\ensuremath{r}},
    description={Tasa de descuento anual utilizada en la evaluación económica},
    sort={05}
}

\newglossaryentry{sym:t}{
    type=symbols,
    name={\ensuremath{t}},
    description={Índice del período (año) del horizonte de evaluación},
    sort={06}
}

\newglossaryentry{sym:n}{
    type=symbols,
    name={\ensuremath{n}},
    description={Cantidad de períodos del horizonte de evaluación},
    sort={07}
}

\newglossaryentry{sym:m}{
    type=symbols,
    name={\ensuremath{m}},
    description={Período en el que el flujo acumulado pasa de negativo a positivo},
    sort={08}
}

\newglossaryentry{sym:fa}{
    type=symbols,
    name={\ensuremath{FA_{m-1}}},
    description={Flujo de fondos acumulado al cierre del último período negativo},
    sort={09}
}

% --- duración de la sesión y precio del resumen automático ---
\newglossaryentry{sym:dmedia}{
    type=symbols,
    name={\ensuremath{\bar{d}}},
    description={Duración media ponderada de una sesión, en minutos},
    sort={10}
}

\newglossaryentry{sym:d}{
    type=symbols,
    name={\ensuremath{d}},
    description={Duración de una clase, en minutos},
    sort={11}
}

\newglossaryentry{sym:presumen}{
    type=symbols,
    name={\ensuremath{P_{\text{resumen}}(d)}},
    description={Precio del resumen automático de una clase de duración \ensuremath{d}},
    sort={12}
}

\newglossaryentry{sym:cia}{
    type=symbols,
    name={\ensuremath{C_{\text{IA}}(d)}},
    description={Costo de inteligencia artificial del resumen de una clase de duración \ensuremath{d}},
    sort={13}
}

\newglossaryentry{sym:ct}{
    type=symbols,
    name={\ensuremath{c_t}},
    description={Costo de transcripción por minuto de audio},
    sort={14}
}

\newglossaryentry{sym:ce}{
    type=symbols,
    name={\ensuremath{c_e}},
    description={Costo por token de entrada del modelo de lenguaje},
    sort={15}
}

\newglossaryentry{sym:tau}{
    type=symbols,
    name={\ensuremath{\tau}},
    description={Tokens de transcripción generados por minuto de clase},
    sort={16}
}

\newglossaryentry{sym:cs}{
    type=symbols,
    name={\ensuremath{c_s}},
    description={Costo por token de salida del modelo de lenguaje},
    sort={17}
}

\newglossaryentry{sym:sigma}{
    type=symbols,
    name={\ensuremath{\sigma}},
    description={Tokens de salida y de instrucciones fijas por resumen},
    sort={18}
}
